\documentclass[../proyecto.tex]{subfiles}
\graphicspath{{\subfix{../images/}}}
\begin{document}

\chapter{Hipótesis y pregunta de investigación}

\section{Hipótesis}

La sintesíntesis, un estilo de vida que consiste en sintetizar sintetizadores, es una práctica de creación artística de herramientas electrónicas y computacionales, que cuando es realizada con los valores de ser de fuente abierta y reparables, genera nuevas autonomías electrónicas.

La sintesíntesis, un estilo de vida que consiste en sintetizar sintetizadores, es una insistencia en la luthería de instrumentos electrónicos y computacionales para emitir ruidos. Su realización de forma popular, entendida como pública, abierta, artesanal y en alternativa a las industrias, emergen nuevas estrategias de diseño y fabricación de dispositivos, aparatos, infraestructuras y metáforas para la producción ruidística de forma modular y distribuida.

\section{Pregunta de investigación}

¿Cómo el diseño y la construcción de sintetizadores modulares de fuente abierta, reparables, y modificables impacta en el acceso al ruido y al desarrollo de formas de conocimiento técnico y autonomía electrónica y computacional latinoamericana?

La sintesíntesis, un estilo de vida que consiste en sintetizar sintetizadores, es una insistencia en la luthería de instrumentos electrónicos y computacionales para emitir ruidos. Al realizarlas de forma popular, entendida como artesanal y en contraposición a la industria, emergen nuevas estrategias de diseño y fabricación de dispositivos, infraestructuras y metáforas para acceder al ruido, de carácter modular, y de fuente abierta. 

% máximo media página

\end{document}